\section{Problem Statement}

Design and implement a microgrid simulator that models renewable energy sources (solar panels, wind turbines), diesel generators, battery storage, and electrical loads. The simulation runs in 15-minute time steps over multiple days. At each step, the system must balance generation with demand by dispatching available sources according to a priority order: renewables first, then battery discharge, then diesel as a last resort. The goal is to minimize diesel usage and battery degradation while meeting all load demands.

\section{Core Concepts}

\begin{itemize}
  \item \textbf{Generation Profile:} Time-series data for each source. Solar follows a bell curve peaking at noon; wind is stochastic with gusts and lulls.
  \item \textbf{Load Profile:} Demand curves per load type. Residential loads peak in morning and evening; commercial loads peak during daytime hours.
  \item \textbf{State of Charge (SoC):} Battery level as a percentage. Operational limits require SoC to stay between 20\% (minimum) and 95\% (maximum) to preserve battery health.
  \item \textbf{Energy Balance:} At every time step: $\text{Generation} + \text{Battery\_Discharge} = \text{Load} + \text{Battery\_Charge} + \text{Losses}$.
  \item \textbf{Dispatch Priority:} Renewables are used first (free, zero emissions), then battery discharge, then diesel generators as a last resort.
\end{itemize}

\section{Key Challenges}

\begin{itemize}
  \item \textbf{Forecasting:} Predict solar irradiance, wind speed, and load demand for upcoming time steps to make better dispatch decisions.
  \item \textbf{Storage Arbitrage:} Store energy when renewable generation is abundant and use it when generation is scarce or demand is high.
  \item \textbf{Battery Degradation:} Penalize deep discharge cycles. Frequent deep cycling reduces battery lifespan and should be avoided.
  \item \textbf{Diesel Minimization:} Track total diesel consumption and runtime. Diesel should only run when renewables and battery cannot meet demand.
  \item \textbf{Multi-Day Patterns:} Handle differences between weekday and weekend load profiles, and varying weather patterns across days.
\end{itemize}

\section{Sample Input}

\begin{lstlisting}[style=json]
{
  "sources": [
    {"id": "solar_1", "type": "solar", "peak_kw": 50},
    {"id": "wind_1", "type": "wind", "peak_kw": 30},
    {"id": "diesel_1", "type": "diesel", "peak_kw": 100}
  ],
  "storage": [
    {
      "id": "battery_1",
      "capacity_kwh": 200,
      "soc": 0.80,
      "charge_rate": 25,
      "discharge_rate": 25
    }
  ],
  "loads": [
    {"id": "residential_1", "base_kw": 40},
    {"id": "commercial_1", "base_kw": 60}
  ],
  "time_series": {
    "solar_profile": [0.0, 0.0, 0.0, 0.0, 0.1, 0.3, 0.6, 0.85, 1.0,
                      0.85, 0.6, 0.3, 0.1, 0.0, 0.0, 0.0],
    "load_profile": [0.4, 0.3, 0.3, 0.3, 0.5, 0.7, 0.9, 1.0, 0.95,
                     0.85, 0.7, 0.6, 0.5, 0.6, 0.8, 1.0]
  },
  "simulation": {
    "time_step_minutes": 15,
    "duration_hours": 72
  }
}
\end{lstlisting}

\vfill
\noindent\rule{\textwidth}{0.4pt}
{\small\textit{For full project requirements, STL restrictions, submission format, and grading criteria, refer to \textbf{base.pdf} (available on the \href{https://github.com/BestSithInEU/cse211-term-projects/releases}{Releases page}).}}
